\section{Trapping Rain Water}
\label{sec:sq-trapping-rain-water}

\problemstatement{We are given $n$ non-negative integers $\textit{height}$
representing an elevation map, where $\textit{height}[i]$ is the height of
a bar of width $1$ at position $i$. Compute how much water is trapped
after raining, assuming water can only be held between bars and runs off
at either open end.}

\begin{patternbox}
\emph{``Water trapped between bars of varying heights''} is a problem
about \textbf{valleys bounded by two walls}: water above any single
position is limited by the \emph{shorter} of the nearest wall to its left
and the nearest wall to its right. Whenever a problem's answer at each
position depends on ``the nearest taller boundary on each side,'' a
monotonic stack -- which is precisely a stack of candidate boundaries,
each taller than the one below it -- is a natural fit, processing the
array once and resolving each ``valley'' as soon as a wall tall enough to
close it off arrives.
\end{patternbox}

\subsection*{Understanding the problem carefully}

Maintain a stack of indices whose heights are non-increasing from bottom
to top (each new index pushed is no taller than the one below it,
otherwise it would immediately start resolving valleys). When we
encounter a bar taller than the index on top of the stack, that means we
have just found a \textbf{right wall} for a valley whose floor is the
popped index, and whose \textbf{left wall} is whatever index remains below
it on the stack. The water trapped above that floor, between the two
walls, is determined by the \emph{shorter} of the two walls (water cannot
rise above the shorter wall without spilling over it) minus the floor's
own height, multiplied by the horizontal distance between the two walls.

\begin{ideabox}[Resolve One Valley Floor per Pop]
Scan left to right. While the current bar is taller than the stack's top:
pop that top index as the \textit{valley floor}. If the stack is now
empty, there is no left wall, so this floor traps no water (stop popping
for this bar). Otherwise, the new stack top is the \textit{left wall}, and
the current index is the \textit{right wall}. The trapped water above the
floor is
\[
  (\min(\textit{height}[\textit{leftWall}],
  \textit{height}[\textit{rightWall}]) - \textit{height}[\textit{floor}])
  \times (\textit{rightWall} - \textit{leftWall} - 1).
\]
Add this to the running total, then continue popping (the same right wall
might close off multiple successively higher floors). Finally, push the
current index.
\end{ideabox}

\subsection*{Why resolving level by level, floor by floor, correctly sums all trapped water}

Every unit of trapped water sits directly above some specific bar (the
``floor'') and below the line connecting the nearest sufficiently tall
walls on either side. By processing each floor exactly once -- the moment
a right wall taller than it is found, with whatever remains on the stack
serving as the left wall -- we account for the entire water column above
that floor in one $O(1)$ computation. Because indices are only ever pushed
once and popped at most once, every floor's contribution is computed
exactly once, and the total accumulated is exactly the total trapped
water, with no double-counting and no gaps.

\subsection*{Algorithm}

\begin{enumerate}
  \item Initialize an empty stack and $\textit{totalWater} \gets 0$.
  \item For $i$ from $0$ to $n-1$:
  \begin{enumerate}
    \item While the stack is not empty and $\textit{height}[i] >
          \textit{height}[\text{stack's top}]$:
    \begin{enumerate}
      \item Pop the top as $\textit{floor}$.
      \item If the stack is now empty: break (no left wall).
      \item $\textit{leftWall} \gets$ the new top; $\textit{width}
            \gets i - \textit{leftWall} - 1$;
            $\textit{boundedHeight} \gets
            \min(\textit{height}[\textit{leftWall}],
            \textit{height}[i]) - \textit{height}[\textit{floor}]$.
      \item $\textit{totalWater} \gets \textit{totalWater} +
            \textit{width} \times \textit{boundedHeight}$.
    \end{enumerate}
    \item Push $i$.
  \end{enumerate}
  \item Return \textit{totalWater}.
\end{enumerate}

\subsection*{Dry run}

\dryrun{Take $\textit{height} = [3,0,2,0,4]$.}

\begin{itemize}
  \item $i=0$ ($h=3$): stack empty. Push $0$. Stack $=[0]$.
  \item $i=1$ ($h=0$): $0 \le \textit{height}[0]=3$, no pop. Push $1$.
        Stack $=[0,1]$.
  \item $i=2$ ($h=2$): $2 > \textit{height}[1]=0$, pop floor $=1$. Stack
        $=[0]$, not empty: leftWall$=0$. Width $=2-0-1=1$.
        boundedHeight $=\min(3,2)-0=2$. Water $+=1\times2=2$.
        $\textit{totalWater}=2$. Now $2 \le \textit{height}[0]=3$, stop
        popping. Push $2$. Stack $=[0,2]$.
  \item $i=3$ ($h=0$): $0\le \textit{height}[2]=2$, no pop. Push $3$.
        Stack $=[0,2,3]$.
  \item $i=4$ ($h=4$): $4>\textit{height}[3]=0$, pop floor$=3$. Left
        wall $=2$. Width $=4-2-1=1$. boundedHeight
        $=\min(2,4)-0=2$. Water $+=1\times2=2$. Total$=4$. Continue:
        $4>\textit{height}[2]=2$, pop floor$=2$. Left wall$=0$. Width
        $=4-0-1=3$. boundedHeight $=\min(3,4)-2=1$. Water
        $+=3\times1=3$. Total$=7$. Now $4>\textit{height}[0]=3$, pop
        floor$=0$. Stack empty -- break (no left wall). Push $4$.
\end{itemize}

Total trapped water: $7$.

\subsection*{C++ implementation}

\begin{lstlisting}
#include <bits/stdc++.h>
using namespace std;

int trapRainWater(vector<int>& height) {
    int n = height.size();
    stack<int> st;
    long long totalWater = 0;

    for (int i = 0; i < n; i++) {
        while (!st.empty() && height[i] > height[st.top()]) {
            int floorIdx = st.top();
            st.pop();
            if (st.empty()) break;

            int leftWall = st.top();
            int width = i - leftWall - 1;
            int boundedHeight = min(height[leftWall], height[i]) - height[floorIdx];
            totalWater += (long long)width * boundedHeight;
        }
        st.push(i);
    }
    return totalWater;
}

int main() {
    vector<int> height = {3, 0, 2, 0, 4};
    cout << "Trapped water: " << trapRainWater(height) << endl;
    return 0;
}
\end{lstlisting}

\begin{complexitybox}
\textbf{Time Complexity.} Each index is pushed once and popped at most
once, giving
\[
  O(n).
\]
\textbf{Space Complexity.} $O(n)$ for the stack in the worst case. (A
two-pointer technique solves this same problem in $O(1)$ extra space, by
tracking running left-max and right-max boundaries directly -- worth
knowing as a further optimization, though the stack-based derivation here
is the one that generalizes most directly to the histogram-area problems
in Sections~\ref{sec:sq-largest-rectangle-histogram}--\ref{sec:sq-maximal-rectangle}.)
\end{complexitybox}

\begin{takeawaybox}
``Water/area bounded by the nearest sufficiently tall walls on both
sides'' is solved by a monotonic stack that resolves one bounded region
per pop, the instant a wall tall enough to close it off arrives. This same
skeleton -- pop a floor, look at what remains as the left boundary, the
trigger as the right boundary, compute a bounded quantity, accumulate --
reappears in Largest Rectangle in Histogram
(Section~\ref{sec:sq-largest-rectangle-histogram}).
\end{takeawaybox}
